% Compilar con pdfLaTeX. Archivo codificado en UTF-8.
\documentclass[11pt,a4paper,oneside,openany]{book}

% Idioma y tipografía.
\usepackage[T1]{fontenc}
\usepackage[utf8]{inputenc}
\usepackage[spanish,es-nodecimaldot,es-noquoting]{babel}
\usepackage{lmodern}
\usepackage{microtype}

% Página y figuras.
\usepackage[margin=2.8cm]{geometry}
\usepackage{graphicx}
\usepackage{xcolor}
\usepackage{booktabs}

% Matemáticas.
\usepackage{amsmath,amssymb,amsthm}
\usepackage{mathtools}
\usepackage{bm}
\usepackage{slashed}

% Entornos numerados por capítulo.
\theoremstyle{plain}
\newtheorem{teorema}{Teorema}[chapter]
\newtheorem{proposicion}[teorema]{Proposición}
\newtheorem{lema}[teorema]{Lema}
\newtheorem{corolario}[teorema]{Corolario}
\theoremstyle{definition}
\newtheorem{definicion}[teorema]{Definición}
\newtheorem{ejemplo}[teorema]{Ejemplo}
\newtheorem{ejercicio}[teorema]{Ejercicio}
\theoremstyle{remark}
\newtheorem*{observacion}{Observación}

% Notación básica; ampliar a medida que se necesite.
\newcommand{\dd}{\mathop{}\!\mathrm{d}}
\newcommand{\ii}{\mathrm{i}}
\newcommand{\ee}{\mathrm{e}}
\newcommand{\Lag}{\mathcal{L}}
\newcommand{\Ham}{\mathcal{H}}
\DeclareMathOperator{\Tr}{Tr}
\DeclareMathOperator{\diag}{diag}
\DeclarePairedDelimiter{\abs}{\lvert}{\rvert}
\DeclarePairedDelimiter{\norm}{\lVert}{\rVert}
\DeclarePairedDelimiter{\bra}{\langle}{\rvert}
\DeclarePairedDelimiter{\ket}{\lvert}{\rangle}
\DeclarePairedDelimiterX{\braket}[2]{\langle}{\rangle}{#1\delimsize\vert #2}

% Enlaces y referencias; cargar al final del preámbulo.
\usepackage[unicode,hidelinks]{hyperref}
\usepackage[spanish,nameinlink,noabbrev]{cleveref}
\hypersetup{pdftitle={Notas de teoría cuántica de campos}}

\title{Notas de teoría cuántica de campos}
\author{}
\date{}

\begin{document}

\chapter{Introducción}

En estas notas vamos a presentar las ideas centrales de la teoría cuántica de campos.

La siguiente estructura es tentativa y podrá ajustarse a medida que avance
la escritura:

\begin{enumerate}
    \item \textbf{Preliminares y convenciones.}
    Repaso de la teoría clásica de campos, simetrías y teorema de Noether.
    Convenciones relativistas, unidades y notación.

    \item \textbf{Formalismo canónico.}
    Cuantización del campo escalar libre, espacio de Fock, operadores de
    creación y destrucción, propagadores y causalidad. Extensión al campo
    de Dirac y a la cuantización fermiónica.

    \item \textbf{Construcción desde partículas, al estilo Weinberg.}
    Representaciones del grupo de Poincaré, estados de una y varias
    partículas, matriz $S$, principio de descomposición en cúmulos y
    construcción de campos locales. Comparación de las hipótesis y los
    resultados de este enfoque con los del formalismo canónico.

    \item \textbf{Interacciones y teoría de perturbaciones.}
    Imagen de interacción, serie de Dyson, teorema de Wick, amplitudes y
    diagramas de Feynman. Funciones de correlación, reducción LSZ y conexión
    con secciones eficaces y tasas de desintegración.

    \item \textbf{Modelos escalares como ejemplos completos.}
    Estudio de la teoría $\lambda\phi^4$ y de un modelo de Yukawa escalar.
    Cálculo de procesos sencillos, desde el lagrangiano hasta los
    observables, y estudio del límite no relativista cuando corresponda.
    Estos modelos servirán también como ejemplos recurrentes en los
    bloques posteriores.

    \item \textbf{Integral de camino.}
    Relación con el formalismo de operadores, funcional generador,
    funciones de correlación y obtención de las reglas de Feynman.

    \item \textbf{Teorías gauge y electrodinámica cuántica (QED).}
    Redundancia gauge, cuantización del campo electromagnético, fijación
    de gauge y acoplamiento a la materia. Identidades de Ward y cálculo de
    procesos sencillos en QED.

    \item \textbf{Renormalización.}
    Divergencias, regularización, contraterminos y condiciones de
    renormalización. Aplicación a modelos escalares y a QED, grupo de
    renormalización y perspectiva de las teorías efectivas. Las primeras
    divergencias podrán introducirse al estudiar los ejemplos escalares.
\end{enumerate}

A lo largo de las notas se procurará hacer explícitas las hipótesis que
sostienen cada resultado importante y su ámbito de validez. Se distinguirán
los resultados exactos de las aproximaciones y los desarrollos
perturbativos, indicando, cuando corresponda, el orden del cálculo y el
régimen de escalas o acoplamientos en que puede utilizarse. Las convenciones
y las hipótesis comunes se declararán una sola vez; las condiciones
adicionales se indicarán junto al resultado que las necesite.

\end{document}
